\chapter{Subgraph matching ingenuo: por qué no escala}
\label{cap:naive-matching}

\begin{intuicion}
El DFS subgraph match del capítulo anterior es \emph{el} baseline.
En grafos chicos o hipótesis con etiquetas selectivas, anda
perfecto. En grafos grandes o hipótesis sin filtros, su complejidad
\(O(\deg^k)\) explota. Este capítulo pone números concretos a esa
explosión: dado un grafo \emph{scale-free} típico (\(\sim 10^7\)
aristas, distribución de grado power-law), una hipótesis de \(k=4\)
sin filtros puede generar trillones de paths. Esa es la motivación
empírica para los \emph{prunes} (próximos dos capítulos), los
\emph{worst-case-optimal joins} (cap.~\ref{cap:lftj}), y los
\emph{semi-join reducers} (cap.~\ref{cap:yannakakis-chain}).
\end{intuicion}

\section{El peor caso teórico}

Recordemos el bound del capítulo~\ref{cap:dfs-baseline}:
\[
\text{paths visitados} \leq |\V_{L_0^V}| \cdot \prod_{i=0}^{k-1} \deg^{+}_{L_i^E, L_{i+1}^V}.
\]

\noindent
Si todas las etiquetas son ``cualquiera'' (sin filtro), \(\deg^{+}\)
es el grado promedio del grafo, \(\bar{d}\). Si todas las
etiquetas son selectivas, cada \(\deg^{+}_{L^E, L^V}\) es una
fracción chica de \(\bar{d}\). El producto \(\prod \deg^{+}\) tiene
crecimiento geométrico en \(k\); con \(\bar{d} = 100\) y \(k = 4\),
eso son \(10^8\) paths \emph{por start}. Multiplicado por
\(|\V_{L_0^V}| = 10^5\) starts: \(10^{13}\) paths totales. A 1~ns
por path (optimistic), 10\,000 segundos = 2.7 horas.

\subsection{El gráfico que vale más que el álgebra}

\begin{figure}[H]
\centering
\begin{tikzpicture}[
    every node/.style={font=\footnotesize}
]
\draw[->, thick] (0,0) -- (8.5,0) node[right]{\(k\)};
\draw[->, thick] (0,0) -- (0,5.5) node[above, align=center]{paths visitados\\\scriptsize (log scale)};

\foreach \x/\lbl in {1/1, 2/2, 3/3, 4/4, 5/5, 6/6, 7/7, 8/8} {
    \draw (\x,-0.05) -- (\x,0.05) node[below=2pt, font=\scriptsize]{\lbl};
}
\foreach \y/\lbl in {1/$10^2$, 2/$10^4$, 3/$10^6$, 4/$10^8$, 5/$10^{10}$} {
    \draw (-0.05,\y) -- (0.05,\y) node[left=2pt, font=\scriptsize]{\lbl};
}

\draw[accent, thick] plot[smooth] coordinates {(1,1.5) (2,2.0) (3,2.4) (4,2.8) (5,3.1) (6,3.4) (7,3.7) (8,4.0)};
\node[accent] at (5,3.6) {\scriptsize típico (con podas)};

\draw[red, thick, dashed] plot[smooth] coordinates {(1,1.5) (2,2.5) (3,3.5) (4,4.5) (5,5.5)};
\node[red] at (5,5.5) [above]{\scriptsize naive $O(\bar{d}^k)$};

\draw[gray, thick, dotted] (0, 4) -- (8, 4);
\node[gray, font=\scriptsize] at (8, 3.85) [left]{\scriptsize cap=10\,000 = $10^4$};
\end{tikzpicture}
\caption{Crecimiento de paths visitados por hunt vs longitud de la
hipótesis. Naive es \(O(\bar{d}^k)\) ---explosión geométrica que
satura el cap a \(k = 3\)---. Con podas, el crecimiento es
sub-exponencial; en grafos típicos los hunts de \(k = 6\) terminan
bajo el cap.}
\label{fig:naive-explosion}
\end{figure}

\section{¿Por qué scale-free empeora todo?}

Las redes corporativas reales NO son grafos aleatorios uniformes.
Son \emph{scale-free}: la distribución de grado sigue una
\emph{power-law},
\(\Pr[\deg(v) = d] \propto d^{-\alpha}\) con \(\alpha \in [2, 3]\).
Eso significa:

\begin{itemize}
    \item La mayoría de los vértices tienen grado bajo (1--10).
    \item Una minoría tiene grado mediano (50--500).
    \item Existe un puñado de \emph{hubs} con grado masivo (\(10^4\)
    a \(10^6\)).
\end{itemize}

\noindent
El DFS naive sufre desproporcionalmente por los hubs:

\begin{itemize}
    \item Un \emph{Domain Controller} con \(10^4\) usuarios
    autenticando contra él tiene grado entrante 10\,000.
    \item Un hunt de la forma
    \texttt{User -[Auth]-> Host -[Connect]-> IP} pasa por el DC en
    el segundo paso. Iterar sus \(10^4\) salidas es \emph{una sola}
    iteración del DFS \emph{en ese vértice}; pero el DFS volverá al
    DC desde \emph{cada uno} de los starts que lo alcanzan.
    \item Si \(10^3\) usuarios alcanzan al DC, el DFS itera
    \(10^4\) salidas \(10^3\) veces \(\Rightarrow\) \(10^7\)
    iteraciones sólo en ese paso.
\end{itemize}

\subsection{El \emph{output} también explota}

Más sutil que la \emph{exploración}: el resultado \emph{en sí}
puede ser enorme. Si un hub tiene grado \(10^4\) y aparece en el
medio de una hipótesis de \(k = 3\), las paths que pasan por él se
multiplican: una arista \emph{entrante} al hub combinada con
\(10^4\) salientes del hub, da \(10^4\) paths \emph{por arista
entrante}. Para \(10^3\) entrantes: \(10^7\) paths.

GraphHunter \emph{cap-eaba} la salida en 10\,000 resultados; pero
el DFS necesita \emph{generarlos para verificarlos}. Aún con cap,
el trabajo del DFS es proporcional al output total, no al cap.
Esto es exactamente el caso ``cap-saturado'' del capítulo
\ref{cap:hunt-result}: el DFS produce más resultados que el cap, y
el costo del emit (no de la búsqueda) es lo que mata.

\section{Tres caminos para escalar}

Cuando el DFS naive explota, hay tres familias de soluciones, cada
una atacando un aspecto distinto:

\subsection{Camino 1: Reducir \emph{starts} con prunes}

Si una fracción significativa de los starts no genera ningún
resultado, descartarlos antes de iniciar el DFS ahorra trabajo.
Dos prunes:

\begin{itemize}
    \item \textbf{NLF prune}: si el primer paso requiere
    \(\geq r\) vecinos del tipo \(L\), descartar \(v\) si
    \(\NLF(v, L) < r\). Local, \(O(|\V|)\).
    \item \textbf{k-hop reachability prune}: si la hipótesis exige
    llegar al tipo \(L_k^V\), descartar \(v\) si \(L_k^V \notin
    \mathrm{bits}[v]\). Transitivo, requiere índice precomputado.
\end{itemize}

Capítulos \ref{cap:prune-nlf} y \ref{cap:prune-khop} desarrollan.

\subsection{Camino 2: Cambiar el algoritmo (WCOJ)}

Para hipótesis con \emph{shape} estructural especial (\(K_4\)), el
DFS naive es estrictamente subóptimo en peor caso: visita
\(O(\deg^3)\) cuando un \emph{worst-case-optimal join} (LFTJ)
visita \(O(m)\). Capítulos \ref{cap:wcoj} y \ref{cap:lftj}.

\subsection{Camino 3: Reducir antes de enumerar (Yannakakis)}

Para hipótesis \(\alpha\)-acíclicas (cadena, estrella, árbol), una
\emph{semi-join reducer} elimina del grafo todas las aristas que
no participan de \emph{ningún} match completo, ANTES de empezar a
enumerar. La enumeración después es trivial. Capítulos
\ref{cap:hipergrafos} y \ref{cap:yannakakis-chain}.

\section{¿Cuándo el DFS naive es suficiente?}

A pesar de su pesadilla teórica, el DFS naive (con podas) anda
perfecto en muchos casos:

\begin{itemize}
    \item \textbf{Hipótesis con etiquetas muy selectivas.} Por
    ejemplo, \texttt{User -[MfaPush]-> Device}. Si MfaPush es raro
    (pocas aristas), el producto \(\deg^{+}_{L^E,L^V}\) es chico.
    \item \textbf{Hipótesis cortas.} \(k = 1\) o \(k = 2\) es DFS
    chico por construcción.
    \item \textbf{Grafos \emph{no} scale-free.} Algunos datasets
    sintéticos o redes muy estructuradas tienen distribuciones de
    grado uniforme; ahí el DFS no sufre los hubs.
\end{itemize}

\noindent
La estrategia del motor: usar el DFS por defecto, ofrecer
operadores especializados cuando la hipótesis tiene la estructura
adecuada. El \emph{planner} (capítulo~\ref{cap:planner}) hace la
elección.

\section{Resumen}

\begin{itemize}
    \item DFS naive: \(O(\bar{d}^k)\) en peor caso. Geométrico en
    \(k\); explota a \(k = 3, 4\) sin podas.
    \item Grafos scale-free amplifican: hubs causan multiplicación
    de paths intermedios.
    \item Tres caminos de escalado: prunes (reducir starts),
    WCOJ (mejor algoritmo), Yannakakis (reducir antes de enumerar).
    \item DFS bien-podado sigue siendo el baseline; los operadores
    especializados se dispatchan cuando la estructura lo permite.
\end{itemize}

\begin{ejercicios}
\begin{enumerate}
    \item \starred{} Para un grafo con \(\bar{d} = 50\) y una
    hipótesis de \(k = 4\) sin filtros sobre \(|\V_{L_0^V}| = 10^4\)
    starts, ¿cuántos paths visita el DFS naive en peor caso?

    \item ¿Cómo cambia el orden de magnitud si las etiquetas filtran
    al 10\% del grado? ¿Y si filtran al 1\%?

    \item \starred{} En un grafo scale-free, los hubs son raros
    pero contribuyen mucho. Modelá: si el grafo tiene \(10^6\)
    vértices con \(\bar{d} = 5\) pero un solo hub con
    \(\deg = 10^5\), ¿cuál es el costo extra del DFS al pasar por
    el hub vs por un vértice promedio?

    \item Diseñá una hipótesis ``imposible'' (que el DFS no termine
    en tiempo razonable). ¿Qué propiedad estructural la convierte
    en imposible? ¿Cómo la detectaría el motor para rechazarla?

    \item ¿Por qué el cap del cuello de botella es el output, no
    la exploración? Pensá: si el DFS encuentra \(10^7\) paths
    válidos pero sólo emite los primeros \(10^4\), ¿el costo de
    \emph{emit} es proporcional al cap o al total?
\end{enumerate}
\end{ejercicios}
